\documentclass[a4paper,12pt]{article}

\usepackage[margin=1in]{geometry}
\usepackage{graphicx}
\usepackage{caption}
\usepackage{subcaption}
\usepackage{hyperref}

\title{Weekly Report 8: Testing Visual Hull Carving with Different Angle Images}

\author{Project Number: 18}

\begin{document}

\date{}

\maketitle

\section{Introduction}

This week, we tested our visual hull carving implementation using different angle pictures from our diamond dataset, especially focusing on larger images.  

After establishing the optimal voxel grid size ($64^3$) in the previous week, we wanted to verify that our carving algorithm works correctly with real silhouette images captured at various rotation angles.

\section{Methodology}

We tested the visual hull carving process with images from different angles in our dataset:

\begin{enumerate}
    \item Loaded diamond images from various rotation angles.
    \item Processed each image to extract clean silhouettes.
    \item Applied visual hull carving to see how different angles contribute to the 3D reconstruction.
    \item Compared results from using different subsets of images.
\end{enumerate}

\subsection{Testing with Different Angles}

We tested our visual hull carving with images captured at different rotation intervals: \textbf{0.9°}, \textbf{1.8°}, and \textbf{3.6°}.  

These different angle intervals allowed us to see how the spacing between images affects the quality of the 3D reconstruction.

\textbf{Implementation in Python:}

\begin{verbatim}
import cv2
import numpy as np

# Test with different angle intervals
angle_intervals = [0.9, 1.8, 3.6]
results = {}

for interval in angle_intervals:
    silhouettes = []
    num_images = int(360 / interval)
    
    for i in range(num_images):
        angle = i * interval
        img_path = f"diamond_{angle:05.1f}.bmp"
        img = cv2.imread(img_path, cv2.IMREAD_GRAYSCALE)
        _, binary = cv2.threshold(img, 127, 255, cv2.THRESH_BINARY)
        silhouettes.append(binary > 127)
    
    # Apply carving with these silhouettes
    carved_model = carve_visual_hull(voxel_grid, silhouettes)
    results[interval] = carved_model
\end{verbatim}

\subsection{Testing with Larger Images}

We also tested with larger resolution images to see if they provide better detail in the reconstructed model.  

Some images in our dataset are larger than the standard 512×512 size, and we wanted to understand how image size affects the carving quality.

\section{Results}

After testing with different angle images, we observed the following:

\begin{table}[h!]
\centering
\caption{Testing Results with Different Angle Intervals}
\begin{tabular}{|c|c|c|c|}
\hline
\textbf{Angle Interval} & \textbf{Number of Images} & \textbf{Result} & \textbf{Processing Time} \\
\hline
3.6° & 100 images & Basic shape, some detail loss & ~6 seconds \\
1.8° & 200 images & Good detail, smooth edges & ~12 seconds \\
0.9° & 400 images & Excellent detail, very smooth & ~24 seconds \\
\hline
\end{tabular}
\label{tab:testing}
\end{table}

\textbf{Key observations:}

\begin{itemize}
    \item Using 3.6° intervals (100 images) gives a basic diamond shape but some fine details are lost between angles.
    \item Using 1.8° intervals (200 images) provides good detail and captures most of the diamond's faceted structure.
    \item Using 0.9° intervals (400 images) gives the best results with excellent detail and very smooth surfaces, but requires more processing time.
    \item Larger resolution images provide better edge definition and finer details in the reconstructed model.
    \item The smaller the angle interval, the more accurate the 3D model becomes, but processing time also increases proportionally.
\end{itemize}

\section{Discussion}

Testing with different angle intervals (0.9°, 1.8°, and 3.6°) helped us understand:

\begin{itemize}
    \item \textbf{Angle interval matters} — smaller intervals (0.9°) provide better coverage and more accurate 3D reconstruction. Larger intervals (3.6°) result in a blockier, less detailed model.
    \item \textbf{Image size affects quality} — larger images capture more detail, which translates to better carving results, especially for fine features like the diamond's facets.
    \item \textbf{Processing time} increases with the number of images, but the quality improvement is worth it. The 1.8° interval provides a good balance between quality and processing time.
    \item \textbf{Optimal interval} — for our dataset, 1.8° intervals (200 images) provide sufficient detail without excessive processing time, making it the most practical choice.
\end{itemize}

The tests confirmed that our visual hull carving algorithm works correctly with real diamond images and that using more angles and larger images produces better results.

\section{Conclusion}

This week, we successfully tested our visual hull carving implementation with different angle pictures from our dataset.  

The key findings include:

\begin{itemize}
    \item The carving algorithm works correctly with real diamond silhouette images.
    \item Smaller angle intervals (0.9°) provide the best quality, while 1.8° offers a good balance.
    \item Larger resolution images provide better detail in the reconstructed 3D model.
    \item The 1.8° interval (200 images) is optimal for our workflow, providing good quality with reasonable processing time.
\end{itemize}

These tests validated our approach and prepared us for the next step: creating a complete pipeline from image loading to 3D model export.

\end{document}
